\begin{abstract}
% A notable byproduct of large language model (LLM) alignment training is the phenomenon of mode collapse: a progressive reduction in output diversity over the course of training that narrows the model’s expressive capacity at inference time. This diminishing in generative diversity is particularly harmful to applications that depend on creativity, open-ended exploration, and pluralistic perspectives.  
% We present \textbf{\method}~(Mode-conditioned Diversity Alignment), a post-training RL algorithm that jointly optimizes quality and diversity of LLM generations. \method~trains a single shared policy across abstract numbered mode tokens to adopt distinct mode-conditioned roles, steering it to seek distinct regions of the high-quality output space without hand-crafted personas or architectural changes. We design a Quality-Gated Diversity Reward (\textbf{\qgdr}) to ensure diversity is rewarded only when outputs meet a prompt-adaptive quality threshold calibrated against a frozen reference policy, usually the initial policy. The gating mechanism ensures that diversity is rewarded only to useful responses, preventing degenerate reward-hacking behaviors such as language switching, unnecessary verbosity and irrelevant responses, while encouraging generation quality. To systematically study quality-diversity tradeoffs, we evaluated our method on a comprehensive evaluation suite composing of a 7-task general capability retention test, covering standard reasoning and instruction-following benchmarks, and a 4-task domain-specific diversity assessment spanning scientific ideation and creative writing. \method outperforms the base model on domain-specific applications by X\% across diversity metrics, while simultaneously improving performance on general capability tasks by Y\%. By jointly enhancing generative quality and diversity, \method~may serve as a drop-in replacement for standard post-training algorithms that otherwise contract the model's exploration space.

% \liweiaddressed{Could we add a github link? It's ok to leave it "under construction" for a bit longer and keep polishing it while making the draft public}

% \liweiaddressed{I think it's better to only mention types of tasks that we do show improvements in the abstract. Otherwise, the reader might expect us to show improvement over all tasks we mention here. Also, I think it would be better overall to remove the math-reasoning mentioning, since people’s immediate reaction might be that math is primarily about logical reasoning, making it less obvious why diversity is relevant there.}.

% \liweiaddressed{It feels we are putting too much emphasis of quality gating at the abstract by spending two long sentences here. It almost occupies 5 lines of text, which takes a lot of attention. But it should just be one part of the overall algorithm. Could we reduce the above two sentences to be 1, and to be no more than 2.5 lines. We can still expand in the intro, but it's a bit disproportional in the abstract.} %while preserving generation quality.

% \liweiaddressed{"We further introduce" makes it sound like we are introducing some other method, yet here it's a key component of the MODA method. Maybe change it to "A key components of MODA is"}

% \liweiaddressed{We can briefly mention MARL?}


A notable byproduct of LLM alignment training is \textit{mode collapse}: the progressive loss of output diversity that narrows a model's expressivity at inference time. This degradation is especially limiting for applications requiring open-ended exploration and pluralistic perspectives, such as scientific ideation and creative writing. We present \textbf{\methodcolor}~(\textbf{Mo}de-conditioned \textbf{D}iversity \textbf{A}lignment), an online algorithm that jointly optimizes generation quality and diversity, inspired by the coordination perspective in multi-agent reinforcement learning (MARL). \method trains a single shared LLM policy conditioned on abstract numbered mode tokens, where each mode acts as an agent that competes to produce outputs distinct from the others. This formulation encourages mode-conditioned agents to explore complementary regions of the high-quality output space by employing a prompt-adaptive quality gating mechanism that calibrates quality thresholds against a frozen reference policy.
We evaluate \method{} on a comprehensive suite of seven general capability tasks and four domain-specific diversity tasks in scientific ideation and creative writing. \method{} improves SBERT diversity by \textbf{265\%} on the Infinite-Chat held-out prompts, while increasing general capability pass@1 by \textbf{10.3\%} over the Qwen3-8B baseline. Compared with the strongest DivPO baseline, \method{} improves SBERT diversity from 0.274 to 0.482 (\textbf{+75.9\%}) and E-Vendi from 2.86 to 4.4 (\textbf{+53.8\%}), while improving average general capability pass@1 by \textbf{7.0\%}. Overall, \method{} provides a drop-in alternative to standard post-training methods that preserves and expands the model's expressive output space while improving quality.


% \method trains a single shared LLM policy conditioned on abstract numbered mode tokens, where each mode acts as an agent that competes to produce outputs distinct from the others. This MARL-inspired formulation encourages mode-conditioned agents to explore complementary regions of the high-quality output space without requiring hand-crafted personas or architectural modifications. To prevent reward hacking, MODA employs a prompt-adaptive quality gating mechanism that calibrates quality thresholds against a frozen reference policy and grants diversity rewards only to responses that satisfy them, mitigating degenerate behaviors such as language switching and verbosity.


% \liwei{This motivating sentence feels a bit too spicy and too abstract for people who are not familiar with mode collapse of LLMs from the beginning. We can make it a bit less "impact"-based but more "evidence"-based (can still mention impact, but after providing evidence): LLM alignment training can induce mode collapse, diminishing output diversity during training and thereby limiting the model’s expressive capacity at inference time. This reduction in diversity may constrain applications that depend on creativity, exploration, and a plurality of perspectives. Starting by saying "LLMs are homogenizing human thought ..." may prime the reader to think we are trying to mitigate how AI's impact on human thoughts, and provide scientific evidence for that, risk setting the wrong expectation of the readers.} 


% \liwei{I think we should soften the framing of the benchmark contribution because we are not introducing new benchmarks for both the general capability and the diversity part. We can say that "through thorough evaluation over standard general capability and reasoning benchmarks, and show there's no capability decrease." For the diversity part, we can frame is as a unified benchmark for measuring model diversity by repurposing previous individual benchmarks, covering xx, xx, xx. We don't want to be come across as overclaiming on benchmark contribution}

% \liwei{Add a elevating sentence to highlight the impact of our method, e.g., \method serves a drop-in replacement of standard post-training algorithm that xx.}

\end{abstract}